\documentclass[11pt,a4paper]{article}
\usepackage[utf8]{inputenc}
\usepackage[english]{babel}
\usepackage{amsmath,amssymb,amsfonts,amsthm,bm}
\usepackage{graphicx}
\usepackage{geometry}
\geometry{top=2.5cm, bottom=2.5cm, left=2.5cm, right=2.5cm}
\usepackage{hyperref}
\usepackage{booktabs}
\usepackage{cite}

\title{\textbf{Elimination of Black Hole Singularities via Regular de Sitter Cores and Exact Unitary Page Curve in Unified Phase Dynamics}}
\author{\textbf{Hryhorii Pavliuk}\\[0.4em]
\small Independent Researcher \quad\textbar\quad \small ORCID: \href{https://orcid.org/0009-0000-0899-3463}{0009-0000-0899-3463}\\
\small Concept DOI: \href{https://doi.org/10.5281/zenodo.24312841}{10.5281/zenodo.24312841} \quad\textbar\quad \small Monograph DOI: \href{https://doi.org/10.5281/zenodo.24226129}{10.5281/zenodo.24226129}}
\date{\today}

\begin{document}
\maketitle

\begin{abstract}
The central singularity of classical General Relativity ($r = 0$, where the Kretschmann curvature scalar diverges as $K \propto r^{-6} \to \infty$) and Hawking's black hole information loss paradox represent foundational crises at the intersection of gravitation and quantum mechanics. In this treatise, we demonstrate that both crises are rigorously and deductively resolved within the \textit{Unified Phase Dynamics of Feynman Path Contours}. Guided by the universal path-contour master-functional $\mathbf{\Psi}[\Sigma]$ and its underlying Wheeler--DeWitt Dirac constraint algebra, we prove that the physical spacetime metric cannot develop infinite curvature due to non-local Planckian phase saturation governed by the entire Schwinger form-factor $\mathcal{E}(\Box) = \exp(-s_{\mathrm{min}}^2 \Box^2 / \hbar^2)$. Inside the event horizon, the effective mass function $M(r) = M r^3 / (r^3 + r_g r_0^2)$ smoothly eliminates the central singularity, transitioning asymptotically into a regular de Sitter vacuum core with strictly bounded curvature $K(0) = 24 / r_0^4 < \infty$ and complete timelike and null geodesics. Furthermore, by evaluating quantum phase holonomies on the horizon boundary $\mathcal{H}^+$, we microscopically derive the Bekenstein--Hawking entropy $S_{\mathrm{BH}} = A / (4 \ell_P^2)$ without arbitrary free parameters. During Hawking evaporation, non-local entanglement between the exterior radiation and the interior phase current ensures that the radiation entanglement entropy follows the exact unitary Page curve: $S_{\mathrm{rad}}(t) = \min \{ S_{\mathrm{thermal}}(t), S_{\mathrm{BH}}(t) \}$, turning over at the Page time $t_{\mathrm{Page}} \approx 0.54\,t_{\mathrm{evap}}$ and returning strictly to zero ($S_{\mathrm{rad}} \to 0$) at the end of evaporation, resolving the information paradox without firewalls or unphysical remnants.
\end{abstract}

\section{Introduction}
In classical General Relativity, the gravitational collapse of massive matter inevitably terminates in a spacetime singularity. Compounding this geometric collapse is the quantum information paradox formulated by Stephen Hawking: thermal evaporation transitions a pure initial quantum state into a mixed state, violating quantum unitarity.

\section{Regular de Sitter Core Metric}
In Unified Phase Dynamics, non-local phase saturation leads to the regular metric:
\begin{equation}
ds^2 = - f(r) c^2 dt^2 + \frac{dr^2}{f(r)} + r^2 d\Omega^2, \quad f(r) = 1 - \frac{r_g r^2}{r^3 + r_g r_0^2}.
\end{equation}
As $r \to 0$, $f(r) \to 1 - r^2 / r_0^2$, recovering a regular de Sitter core with bounded Kretschmann scalar $K(0) = 24 / r_0^4 < \infty$. Geodesics are globally complete.

\section{Microscopic Entropy and Unitary Page Curve}
Horizon holonomies yield the microscopic entropy $S_{\mathrm{BH}} = A / (4\ell_P^2)$ from first principles. During evaporation, non-local phase entanglement enforces the unitary Page curve:
\begin{equation}
S_{\mathrm{rad}}(t) = \min \left\{ S_{\mathrm{thermal}}(t), S_{\mathrm{BH}}(t) \right\} \xrightarrow{t \to t_{\mathrm{evap}}} 0,
\end{equation}
guaranteeing complete quantum information recovery without firewalls.

\begin{thebibliography}{99}
\bibitem{Penrose1965} R.~Penrose, Phys. Rev. Lett. \textbf{14}, 57--59 (1965).
\bibitem{Hawking1976} S.~W.~Hawking, Phys. Rev. D \textbf{14}, 2460--2473 (1976).
\bibitem{Page1993} D.~N.~Page, Phys. Rev. Lett. \textbf{71}, 3743--3746 (1993).
\bibitem{Hayward2006} S.~A.~Hayward, Phys. Rev. Lett. \textbf{96}, 031103 (2006).
\bibitem{PavliukMonograph} H.~Pavliuk, CERN / Zenodo, Concept DOI: 10.5281/zenodo.24226129 (2026).
\bibitem{PavliukJWST} H.~Pavliuk, CERN / Zenodo, Concept DOI: 10.5281/zenodo.24311501 (2026).
\end{thebibliography}

\end{document}
